\documentclass[a4paper,12pt]{article}

\usepackage[T2A]{fontenc}
\usepackage[utf8]{inputenc}
\usepackage[russian]{babel}

\usepackage{amsmath, amssymb, amsthm, mathtools}
\usepackage{helvet}
\renewcommand{\familydefault}{\sfdefault}
\usepackage{icomma}
\usepackage{geometry}
\usepackage{microtype}
\usepackage{tikz}
\usepackage{tkz-euclide}
\usepackage{pgfplots}
\pgfplotsset{compat=1.18}
\usepackage{tabularx, booktabs, longtable}
\usepackage{enumitem}
\usepackage{hyperref}
\usepackage{parskip}
\usepackage{fancyhdr}
\usepackage{setspace}
\usepackage{xcolor}

\geometry{left=18mm,right=18mm,top=18mm,bottom=20mm}

\definecolor{accent}{HTML}{008080}
\definecolor{darkaccent}{HTML}{004D4D}
\definecolor{textgray}{HTML}{333333}
\definecolor{softgray}{HTML}{F2F6F6}
\definecolor{linegray}{HTML}{B8C8C8}

\hypersetup{hidelinks}
\onehalfspacing
\setlength{\parindent}{0pt}
\setlength{\parskip}{6pt}

\pagestyle{fancy}
\fancyhf{}
\lhead{\textcolor{darkaccent}{Показательные неравенства}}
\rhead{\textcolor{darkaccent}{Николь}}
\cfoot{\thepage}
\renewcommand{\headrulewidth}{0.4pt}
\renewcommand{\headrule}{\hbox to\headwidth{\color{linegray}\leaders\hrule height \headrulewidth\hfill}}

\setlist[itemize]{leftmargin=1.2cm,itemsep=2pt,topsep=3pt}
\setlist[enumerate]{leftmargin=1.1cm,itemsep=3pt,topsep=3pt}

\newcommand{\sectionline}{\vspace{2mm}{\color{linegray}\hrule height 0.5pt}\vspace{2mm}}
\newcommand{\important}[1]{\textcolor{darkaccent}{\textbf{#1}}}
\newcommand{\odz}{\text{ОДЗ}}
\newcommand{\R}{\mathbb{R}}

\newcommand{\intervalaxis}[5]{
\begin{center}
\begin{tikzpicture}[scale=1]
  \draw[->,line width=0.7pt] (-0.5,0) -- (7.2,0);
  \node[right] at (7.2,0) {$#1$};
  #2
  #3
  #4
  #5
\end{tikzpicture}
\end{center}
}

\begin{document}

\begin{titlepage}
\thispagestyle{empty}
\begin{center}
\vspace*{28mm}

{\Large \textcolor{darkaccent}{\textbf{Чек-лист}}}

\vspace{6mm}

{\huge \textbf{Показательные неравенства}}\\[2mm]
{\Large \textbf{замены, вынесение, деление на положительное выражение}}

\vspace{18mm}

{\large Лёвин Артём Александрович}\\[1mm]
{\large эксклюзивно для Николь}

\vspace{12mm}

{\large \today}

\vfill

\begin{tikzpicture}[x=1cm,y=1cm]
\draw[accent,line width=1.1pt]
  (-5.7,0) .. controls (-3.8,0.9) and (-2.1,-0.8) .. (-0.3,0.15)
  .. controls (1.6,1.05) and (3.5,-0.65) .. (5.7,0.2);
\draw[linegray,line width=0.6pt]
  (-5.2,-0.35) .. controls (-3.0,0.25) and (-1.4,-0.5) .. (0.4,-0.2)
  .. controls (2.4,0.2) and (3.7,-0.3) .. (5.2,-0.1);
\end{tikzpicture}

\vspace{12mm}
\end{center}
\end{titlepage}

\setcounter{page}{2}

\section*{\textcolor{darkaccent}{1. Главная идея темы}}

\important{Показательное неравенство} содержит переменную в показателе степени:
\[
a^{f(x)} \quad \text{или} \quad a^{f(x)}\cdot b^{g(x)}.
\]

Цель решения: привести выражения к виду, где можно использовать одно из трёх действий:
\[
\text{замена}, \qquad
\text{вынесение общего множителя}, \qquad
\text{сравнение степеней при одинаковом основании}.
\]

\sectionline

\section*{\textcolor{darkaccent}{2. Базовые правила степеней}}

\begin{center}
\renewcommand{\arraystretch}{1.45}
\begin{tabularx}{\linewidth}{@{}>{\bfseries}lX@{}}
\toprule
Правило & Как используется в решении \\
\midrule
$a^m\cdot a^n=a^{m+n}$ &
Если основания одинаковые, показатели складываются. \\

$\dfrac{a^m}{a^n}=a^{m-n}$ &
Если основания одинаковые, показатели вычитаются. \\

$(a^m)^n=a^{mn}$ &
Степень степени раскрывается умножением показателей. \\

$a^{-n}=\dfrac1{a^n}$ &
Отрицательные степени лучше переводить в дроби. \\

$(ab)^n=a^n b^n$ &
Составные числа удобно раскладывать на простые множители. \\
\bottomrule
\end{tabularx}
\end{center}

\important{Важно.} Составные числа, корни и дроби в показательных задачах почти всегда нужно переводить в степенной вид:
\[
10^x=2^x\cdot5^x,\qquad
25=5^2,\qquad
\frac{8}{27}=\frac{2^3}{3^3}.
\]

\sectionline

\section*{\textcolor{darkaccent}{3. Сравнение степеней с одинаковым основанием}}

Если
\[
a>1,
\]
то функция $a^x$ возрастает, поэтому знак неравенства сохраняется:
\[
a^{f(x)}\le a^{g(x)}
\quad\Longleftrightarrow\quad
f(x)\le g(x).
\]

Если
\[
0<a<1,
\]
то функция $a^x$ убывает, поэтому знак неравенства меняется:
\[
a^{f(x)}\le a^{g(x)}
\quad\Longleftrightarrow\quad
f(x)\ge g(x).
\]

\important{Типичная ошибка:} нельзя просто «убрать основания», не проверив, больше основание единицы или меньше.

\sectionline

\section*{\textcolor{darkaccent}{4. Замена переменной}}

Если в неравенстве многократно встречается одно и то же показательное выражение, удобно сделать замену.

Например:
\[
3^{2x}-4\cdot3^x+3\le0.
\]

Так как
\[
3^{2x}=(3^x)^2,
\]
обозначим:
\[
t=3^x,\qquad t>0.
\]

Тогда:
\[
t^2-4t+3\le0.
\]

Решаем неравенство по $t$:
\[
(t-1)(t-3)\le0.
\]

\begin{center}
\begin{tikzpicture}[scale=1]
\draw[->,line width=0.7pt] (-0.4,0)--(6.4,0);
\node[right] at (6.4,0) {$t$};

\filldraw[fill=white,draw=black,line width=0.8pt] (1.8,0) circle (2pt);
\filldraw[fill=white,draw=black,line width=0.8pt] (4.6,0) circle (2pt);
\node[below] at (1.8,-0.05) {$1$};
\node[below] at (4.6,-0.05) {$3$};

\node[above] at (0.8,0.15) {$+$};
\node[above] at (3.2,0.15) {$-$};
\node[above] at (5.5,0.15) {$+$};

\draw[line width=1pt,accent] (1.8,0.28)--(4.6,0.28);
\node[above,accent] at (3.2,0.35) {$t\in[1;3]$};
\end{tikzpicture}
\end{center}

Возвращаемся к $x$:
\[
1\le t\le3
\quad\Longleftrightarrow\quad
1\le3^x\le3.
\]

Так как
\[
1=3^0,\qquad 3=3^1,
\]
получаем:
\[
3^0\le3^x\le3^1
\quad\Longleftrightarrow\quad
0\le x\le1.
\]

\[
\boxed{x\in[0;1]}.
\]

\important{Главное правило оформления:} после замены сначала решаем неравенство по новой переменной, затем аккуратно делаем обратную замену.

\sectionline

\section*{\textcolor{darkaccent}{5. Метод интервалов после замены}}

После замены часто получается рациональное неравенство. Его решают методом интервалов.

Пример:
\[
\frac{2(4-3^x)}{3^x-1}\ge0.
\]

Делаем замену:
\[
t=3^x,\qquad t>0.
\]

Получаем:
\[
\frac{2(4-t)}{t-1}\ge0.
\]

Критические точки:
\[
t=1 \quad \text{--- корень знаменателя, точка выколота},
\]
\[
t=4 \quad \text{--- корень числителя, точка закрашена}.
\]

\begin{center}
\begin{tikzpicture}[scale=1]
\draw[->,line width=0.7pt] (-0.3,0)--(7,0);
\node[right] at (7,0) {$t$};

\draw[line width=0.7pt,linegray] (0.7,0.35)--(0.7,-0.35);
\node[below] at (0.7,-0.05) {$0$};
\node[above,linegray] at (0.7,0.4) {$t>0$};

\filldraw[fill=white,draw=black,line width=0.8pt] (2.6,0) circle (2pt);
\filldraw[fill=black,draw=black,line width=0.8pt] (5.2,0) circle (2pt);
\node[below] at (2.6,-0.05) {$1$};
\node[below] at (5.2,-0.05) {$4$};

\node[above] at (1.6,0.15) {$-$};
\node[above] at (3.9,0.15) {$+$};
\node[above] at (6.0,0.15) {$-$};

\draw[line width=1pt,accent] (2.68,0.28)--(5.2,0.28);
\node[above,accent] at (3.95,0.36) {$t\in(1;4]$};
\end{tikzpicture}
\end{center}

Возвращаемся к $x$:
\[
1<3^x\le4.
\]

Левая часть:
\[
1<3^x
\quad\Longleftrightarrow\quad
3^0<3^x
\quad\Longleftrightarrow\quad
x>0.
\]

Правая часть:
\[
3^x\le4
\quad\Longleftrightarrow\quad
x\le\log_3 4.
\]

Итог:
\[
\boxed{x\in(0;\log_3 4]}.
\]

\sectionline

\section*{\textcolor{darkaccent}{6. Выколотая точка остаётся выколотой}}

Если одна и та же точка является одновременно корнем числителя и корнем знаменателя, она всё равно выколота.

Пример:
\[
\frac{(t-2)(t-3)}{(t-2)^2(t+2)}\le0.
\]

Точка $t=2$ встречается в числителе и знаменателе, но знаменатель не может быть равен нулю, значит:
\[
t=2 \quad \text{выколота}.
\]

\important{Правило.}
Если точка выколота хотя бы один раз, она выколота всегда.

\sectionline

\section*{\textcolor{darkaccent}{7. Чётная и нечётная кратность}}

При методе интервалов важно учитывать, сколько раз встречается множитель.

\begin{center}
\renewcommand{\arraystretch}{1.45}
\begin{tabularx}{\linewidth}{@{}lX@{}}
\toprule
Ситуация & Что происходит со знаком \\
\midrule
$(x-a)^1$, $(x-a)^3$, $(x-a)^5$ & Нечётная кратность: знак меняется. \\
$(x-a)^2$, $(x-a)^4$, $(x-a)^6$ & Чётная кратность: знак не меняется. \\
\bottomrule
\end{tabularx}
\end{center}

\sectionline

\section*{\textcolor{darkaccent}{8. Деление на выражение с переменной}}

В неравенствах нельзя делить на произвольное выражение с переменной, потому что неизвестно, положительное оно или отрицательное.

Например, делить на $x$ без дополнительных условий нельзя:
\[
x\cdot A(x)>x\cdot B(x).
\]

Но на показательное выражение делить можно:
\[
7^x>0,\qquad 3^{2x}>0,\qquad 5^{\sqrt{x}}>0.
\]

Поэтому, если делим неравенство на $7^x$, знак не меняется:
\[
\frac{A(x)}{7^x}>\frac{B(x)}{7^x}.
\]

\important{Обязательно показывайте, почему деление корректно:}
\[
7^x>0.
\]

\sectionline

\section*{\textcolor{darkaccent}{9. Приём: собрать степени слева, обычные числа справа}}

Пример:
\[
5^{2x}7^{3x}+2\cdot7^{2x}5^{3x}\le\frac32\cdot7^{2x}5^{3x}.
\]

Так как
\[
7^{2x}5^{3x}>0,
\]
делим обе части на это выражение:
\[
\frac{5^{2x}7^{3x}}{7^{2x}5^{3x}}+2\le\frac32.
\]

Сокращаем степени:
\[
\frac{7^x}{5^x}+2\le\frac32.
\]

Тогда:
\[
\left(\frac75\right)^x\le -\frac12.
\]

Левая часть всегда положительна:
\[
\left(\frac75\right)^x>0,
\]
а правая часть отрицательна:
\[
-\frac12<0.
\]

Значит, решений нет:
\[
\boxed{\emptyset}.
\]

\important{Полезная проверка:} всегда оценивайте, какие значения может принимать левая и правая части. Иногда неравенство невозможно уже по знаку.

\sectionline

\section*{\textcolor{darkaccent}{10. Группировка и двойное вынесение}}

Если в неравенстве четыре слагаемых, часто работает группировка.

Пример:
\[
25\cdot2^x-10^x-25+5^x>0.
\]

Представляем:
\[
10^x=2^x\cdot5^x.
\]

Тогда:
\[
25\cdot2^x-2^x5^x-25+5^x>0.
\]

Группируем:
\[
2^x(25-5^x)-(25-5^x)>0.
\]

Выносим общую скобку:
\[
(25-5^x)(2^x-1)>0.
\]

Преобразуем:
\[
(5^2-5^x)(2^x-2^0)>0.
\]

На ОДЗ данное неравенство знака совпадает с:
\[
(5-1)(2-x)(2-1)(x-0)>0.
\]

Положительные множители $5-1$ и $2-1$ на знак не влияют:
\[
(2-x)x>0.
\]

Метод интервалов:
\[
x\in(0;2).
\]

\[
\boxed{x\in(0;2)}.
\]

\important{Ключевая идея:}
если выражение имеет вид
\[
a^m-a^n,
\]
то знак разности совпадает со знаком произведения:
\[
(a-1)(m-n).
\]

\sectionline

\section*{\textcolor{darkaccent}{11. Правило знака для разности степеней}}

Для $a>0$, $a\ne1$:
\[
a^u-a^v
\]
знака совпадает с
\[
(a-1)(u-v).
\]

Поэтому:
\[
a^u-a^v>0
\quad\Longleftrightarrow\quad
(a-1)(u-v)>0.
\]

Примеры:
\[
5^x-5^2>0
\quad\Longleftrightarrow\quad
(5-1)(x-2)>0
\quad\Longleftrightarrow\quad
x>2.
\]

\[
\left(\frac13\right)^x-\left(\frac13\right)^2>0
\quad\Longleftrightarrow\quad
\left(\frac13-1\right)(x-2)>0
\quad\Longleftrightarrow\quad
x<2.
\]

\important{Почему это удобно:} правило автоматически учитывает, возрастает или убывает показательная функция.

\sectionline

\section*{\textcolor{darkaccent}{12. Задачи с корнем в показателе}}

Если в показателе степени есть $\sqrt{x}$, сначала пишем ограничение:
\[
x\ge0.
\]

Пример:
\[
6^{2\sqrt{x}+4}>3^{3\sqrt{x}}\cdot2^{\sqrt{x}+8}.
\]

Обозначим:
\[
t=\sqrt{x},\qquad t\ge0.
\]

Тогда:
\[
6^{2t+4}>3^{3t}\cdot2^{t+8}.
\]

Представим $6$ как $2\cdot3$:
\[
3^{2t+4}2^{2t+4}>3^{3t}2^{t+8}.
\]

Делим на положительное выражение $3^{3t}2^{t+8}>0$:
\[
\frac{3^{2t+4}2^{2t+4}}{3^{3t}2^{t+8}}>1.
\]

Вычитаем показатели:
\[
3^{2t+4-3t}\cdot2^{2t+4-t-8}>1.
\]

Получаем:
\[
3^{4-t}\cdot2^{t-4}>1.
\]

Приводим к одному основанию отношения:
\[
\frac{3^{4-t}}{2^{4-t}}>1
\quad\Longleftrightarrow\quad
\left(\frac32\right)^{4-t}>1.
\]

Так как
\[
\frac32>1,
\]
получаем:
\[
4-t>0,
\qquad t<4.
\]

Возвращаемся к $x$:
\[
\sqrt{x}<4.
\]

С учётом $x\ge0$:
\[
0\le x<16.
\]

\[
\boxed{x\in[0;16)}.
\]

\sectionline

\section*{\textcolor{darkaccent}{13. Оформление решения на экзамене}}

\begin{enumerate}
\item Пишите решение в один столбик, чтобы было видно, что из чего следует.
\item ОДЗ и ограничения записывайте как полноценную часть решения.
\item После замены обязательно указывайте ограничение на новую переменную.
\item Не сокращайте выражения, если можете потерять выколотые точки.
\item Не делите на выражение с переменной без проверки его знака.
\item Не ставьте знаки на числовой прямой там, где выражение не существует.
\item В обратной замене переходите к системе или совокупности, если интервалов несколько.
\item Все составные числа, дроби и корни по возможности переводите в степенной вид.
\end{enumerate}

\sectionline

\section*{\textcolor{darkaccent}{14. Типичные ошибки}}

\begin{center}
\renewcommand{\arraystretch}{1.45}
\begin{tabularx}{\linewidth}{@{}p{0.44\linewidth}X@{}}
\toprule
Ошибка & Как правильно \\
\midrule
Писать $2\cdot3^x=6^x$ &
Так нельзя. Верно: $2\cdot3^x$ остаётся произведением числа $2$ и степени $3^x$. \\

Считать $2^x\cdot3^x$ отдельным несократимым выражением &
Можно объединить: $2^x\cdot3^x=6^x$. \\

Делить неравенство на $x$ без условий &
Нельзя, потому что знак $x$ неизвестен. \\

Делить на $7^x$ и не объяснять почему &
Можно, но лучше написать: $7^x>0$. \\

После замены сразу переходить к $x$ без решения по $t$ &
Сначала решаем неравенство по $t$, затем делаем обратную замену. \\

Забывать условие $t>0$ при $t=a^x$ &
Обязательно: если $t=a^x$, то $t>0$. \\

Ставить знак на промежутке, где выражение не существует &
В запрещённой области знак не ставим. \\
\bottomrule
\end{tabularx}
\end{center}

\newpage

\section*{\textcolor{darkaccent}{Тренировочный блок}}

Решите неравенства. Упражнения расположены по нарастающей сложности.

\begin{enumerate}
\item
\[
3^{2x}-4\cdot3^x+3\le0.
\]

\item
\[
\frac{2(4-3^x)}{3^x-1}\ge0.
\]

\item
\[
5^{x+1}-25\le0.
\]

\item
\[
\frac{(3^x-2)(3^x-3)}{(3^x-2)^2(3^x+2)}\le0.
\]

\item
\[
25\cdot2^x-10^x-25+5^x>0.
\]

\item
\[
7^{x+1}-49\ge0.
\]

\item
\[
5^{2x}7^{3x}+2\cdot7^{2x}5^{3x}\le\frac32\cdot7^{2x}5^{3x}.
\]

\item
\[
6^{2\sqrt{x}+4}>3^{3\sqrt{x}}\cdot2^{\sqrt{x}+8}.
\]

\item
\[
4^{\sqrt{x}-1}\cdot27^{\sqrt{x}+2}\le
2^{3\sqrt{x}-3}\cdot9^{\sqrt{x}+5}.
\]

\item
\[
\frac{3^{2x}-5\cdot3^x+6}{3^{2x}-4}\ge0.
\]
\end{enumerate}

\newpage

\section*{\textcolor{darkaccent}{Ответы}}

\begin{center}
\renewcommand{\arraystretch}{1.55}
\begin{tabularx}{\linewidth}{@{}cX@{}}
\toprule
№ & Ответ \\
\midrule
1 &
\[
x\in[0;1]
\]
\\

2 &
\[
x\in(0;\log_3 4]
\]
\\

3 &
\[
x\le1
\]
\\

4 &
\[
x\in(-\infty;\log_3 2)\cup[1;+\infty)
\]
\\

5 &
\[
x\in(0;2)
\]
\\

6 &
\[
x\ge1
\]
\\

7 &
\[
\emptyset
\]
\\

8 &
\[
x\in[0;16)
\]
\\

9 &
\[
x\in[0;4]
\]
\\

10 &
\[
x\in(-\infty;\log_3 2)\cup[\log_3 3;+\infty)
\]
\\
\bottomrule
\end{tabularx}
\end{center}

\end{document}